\documentclass[11pt]{article}
\usepackage{amsmath,amssymb,amsthm}
\usepackage[margin=1.1in]{geometry}
\newtheorem{theorem}{Theorem}
\newtheorem{conjecture}{Conjecture}
\theoremstyle{remark}
\newtheorem{remark}{Remark}

\title{The multiplication crystal:\\ dimension, diffraction, and the symmetries of arithmetic}
\author{Carl Gribble\thanks{Draft v0.2, expository companion to the research papers
\cite{CompanionI,CompanionII,CompanionIII} (same author). Developed in collaboration with Claude
(Anthropic); every numerical table in this essay was machine-verified as described in Section~10.
Responsibility for the content rests with the author. A full disclosure of AI use appears at the end
of the paper.}}
\date{August 2026 --- draft}

\begin{document}
\maketitle

\begin{abstract}
The positive integers under multiplication form a crystal: a free structure whose axes are the
primes, whose points are exponent vectors, and whose one-dimensional shadow --- the numbers in their
usual order --- is the aperiodic point set that number theory studies. This essay is a guided tour of
that crystal, conducted with a computer. We compute its dimension and find it fractal: the
self-similarity of the multiplicative structure forces a Moran equation whose form is $\zeta(s) = 2$,
with real solution $\rho = 1.728647\ldots$ (Kalm\'ar's exponent), and the total mass of the crystal's
description tower is verified to follow the law $C y^{\rho}$ to a quarter of a percent at $y = 10^6$
--- a law with a thermodynamic reading: $\rho$ is the Hagedorn point of the gas of ordered
factorizations, and its pole residue is the same constant again.
We diffract it and find a reciprocal crystal: waves shone on the primes' logarithmic positions
produce Bragg peaks at the Riemann zeros (verified to $\pm 0.005$ from a sieve to $2\times 10^5$),
while, in the converse experiment of Landau's 1912 theorem, waves built from $300$ zeros ring at the
primes and prime powers of a window and at nothing else; behind both experiments, the Guinand--Weil
explicit formula is paired by machine against Gaussian tests to twenty-three decimals --- the
diffraction of the weighted prime comb as an identity, not a fit, with the honest clause that the
bare shadow is no Delone set stated and priced. Finally we chart its symmetries: maximally
symmetric as a bare crystal (every permutation of the prime axes is an automorphism), perfectly rigid
as arithmetic (the automorphism group of $(\mathbb{N}, +, \times)$ is trivial), with the sizes
$\log 2, \log 3, \log 5, \ldots$ acting as the symmetry-breaking field between the two. The surviving
symmetries are named and, where possible, measured: scale invariance; the congruence point groups and
their $2016$ fine structure (consecutive primes repel in their last digits: $16.3\%$ of pairs against
a naive $25\%$, measured to $10^7$); the hidden Galois symmetry with the primes as its census; the
exact mirror symmetry of the reciprocal crystal (verified to $10^{-22}$); and its unitary statistical
symmetry class (one near-collision among $499$ zero spacings where structureless randomness predicts
a hundred). The symmetry-breaking language is meanwhile a theorem with an address --- Bost and
Connes's phase transition, whose order parameter the machine watches condense through $\beta = 1$
--- and the weld's wildness is the logicians' trichotomy: ruler and crystal separately decidable,
welded arithmetic undecidable already at one tick of the successor. In this language the Riemann
Hypothesis is a confinement statement: the diffraction
pattern, provably symmetric about its mirror, is conjectured to lie on it. A final comparative section
rebuilds every chapter in Weil's mirror universe $\mathbb{F}_q[t]$, where the axis lengths are
quantized, and finds each aperiodic pathology replaced by an exactness --- closed-form prime counts,
integer dimension, a Mertens function that is identically zero, and a Riemann Hypothesis that is a
theorem, machine-verified here on an elliptic curve at $428$ primes; a third universe --- finite
graphs, their zeta functions rational and their Riemann Hypothesis equivalent to the Ramanujan
property --- is checked per instance, in both directions --- isolating the
incommensurability of the lengths $\log p$ as the single source of the integers' difficulty. A closing analysis
localizes the source once more: by Ostrowski's theorem the incommensurability lives at the unique
archimedean place, and, fused with the absence of Frobenius, traces to one dial --- characteristic
zero. A final experimental section
builds a weld laboratory and a skeleton-null scanner (first light: the consecutive-residue repulsion
rediscovered unprompted; zero new candidates), demonstrates a six-stage scan-to-theorem pipeline on
the $\ell$-adic depth law of $p-1$, and names the programme's organizing claim as the two-skeleton
conjecture: the crystal's detectable order is its congruence skeleton and its spectral skeleton, and
nothing else. Nothing in this essay is
new mathematics; every theorem is attributed. What is offered is the tour, its organizing thesis ---
number theory as the physics of a broken crystal symmetry --- and the machine demonstrations, which
anyone can rerun.
\end{abstract}

\section{A plane of ones}

Begin where arithmetic began: with pebbles. A number is a collection of ones, and the Greeks read
multiplication geometrically --- a product of two numbers is a rectangle of ones, a product of three
a box, and Euclid's \emph{Elements} speaks of plane numbers and solid numbers accordingly
\cite{Euclid}. In this picture a prime is a \emph{linear number}: a quantity of ones that can form
only a line, never a proper rectangle. The multiplication table is the catalogue of all rectangles,
and the primes are the row counts missing from its interior.

The research companions to this essay take that picture seriously as a computational instrument. The
first \cite{CompanionI} turns the de-duplication of the multiplication table into an exact,
prime-free formula system for prime counts and prime sums; the second \cite{CompanionII} builds
certified measuring instruments --- generalized factorials --- that weigh the primes in a window
using only the primes below it; the third \cite{CompanionIII} proves that such measurements determine
the primes individually, and prices the determination in bits. This essay is the view from above:
what single object have those papers been probing, and what is its geometry?

\section{The crystal and its shadow}

Unique factorization is a coordinate statement. Write $n = 2^{a_1} 3^{a_2} 5^{a_3} \cdots$ and the
number $n$ becomes the integer vector $(a_1, a_2, a_3, \ldots)$; multiplication of numbers becomes
addition of vectors. The positive integers under multiplication are therefore a \emph{lattice} --- a
crystal --- of infinite rank, one axis per prime, the numbers themselves its points. We will call it
the multiplication crystal. (Formally: a \emph{crystal} is a free commutative monoid equipped with an additive length
functional with locally finite axis spectrum; the multiplication crystal is the instance with axes
the primes and lengths $\log p$. The companion definitions note~\cite{Defs} fixes this vocabulary,
in which every chapter below computes one invariant of that one object.) Two features deserve
immediate notice. First, the primes are not special
points of the crystal; they are its \emph{axes}. Second, the crystal is free: no relations bind the
axes, which is unique factorization restated.

The numbers as we meet them --- ordered by size on a line --- are a one-dimensional shadow of this
crystal. The projection is the logarithm: $\log n = a_1 \log 2 + a_2 \log 3 + \cdots$, a linear
functional on the exponent lattice with weights $\log p$. The image is an aperiodic point set: the
gaps $\log(n{+}1) - \log n$ shrink irregularly, and no finite pattern repeats. But the disorder is an
artifact of projection. The weights $\log 2, \log 3, \log 5, \ldots$ share no common measure --- a
rational relation among them would exhibit a nontrivial product of primes equal to $1$ --- so the
projection throws a perfectly regular infinite-dimensional crystal onto a line at incommensurable
angles. The mechanism --- a periodic structure of high rank projected at incommensurable angles ---
is exactly the cut-and-project mechanism by which the model sets of mathematical quasicrystal
theory are built \cite{BaakeGrimm}. One honesty clause must ride with the analogy, and Section~4
pays it in full: the shadow itself violates that theory's geometric axioms (its gaps shrink, so it
is not uniformly discrete --- not a Delone set, hence not a model set). The crystalline order
survives not in the bare point set but in the \emph{weighted combs} it carries, and it is there
that diffraction becomes a theorem, to which we will return. A
typical point of the shadow at height $n$ engages about $\log\log n$ of the crystal's axes
\cite{HardyRamanujan}; the primes are the points engaging exactly one --- flattest in the crystal,
and, as positions on the line, the least predictable.

\section{The dimension of the tower}

The crystal is exactly self-similar: multiplying every point by $d$ maps the structure into itself,
for every $d \ge 2$. (This one-way family of translations is what powers the recursions of
\cite{CompanionI}, whose regions at scale $y$ decompose into copies at scales $y/d$.) Self-similar
structures have a canonical dimension, computed by Moran's principle: if a set is built from copies
of itself at scales $r_1, r_2, \ldots$, its similarity dimension is the exponent $s$ with
$\sum_i r_i^{\,s} = 1$. For the multiplication crystal the ratios are $1/d$ for every $d \ge 2$, and
Moran's equation becomes
\[
\sum_{d \ge 2} \frac{1}{d^{\,s}} \;=\; 1, \qquad\text{that is,}\qquad \boxed{\;\zeta(s) = 2\;.}
\]
The zeta function, which will reappear below as the crystal's diffraction apparatus, enters first as
its \emph{dimension formula}. The real solution is
\[
\rho \;=\; 1.728647238998\ldots
\]

What does this non-integer dimension measure? Not the shadow, and not a single layer of the
description space of \cite{CompanionI} --- layer $k$ (ordered products of $k$ factors $\ge 2$) has
generating function $(\zeta(s)-1)^k$, whose pole at $s = 1$ has order exactly $k$, an integer
dimension for each layer, verified numerically in our runs to four decimals. The fractal exponent
governs the \emph{whole tower at once}: the total number $G(y)$ of factorization-descriptions of all
numbers up to $y$, every box in every layer, whose generating function is the geometric series
$\sum_k (\zeta - 1)^k = 1/(2 - \zeta(s))$. Stacking the integer-dimensional layers moves the
rightmost singularity from $s = 1$ out to $s = \rho$: the tower's dimension is where the layer series
diverges. This is a classical theorem of Kalm\'ar \cite{Kalmar}: $G(y) \sim C y^{\rho}$ with
$C = -1/(\rho\, \zeta'(\rho)) = 0.3181736\ldots$; the Moran reading --- $\rho$ as the similarity
dimension of the multiplication crystal --- is, we believe, the right way to see \emph{why} the
equation is $\zeta(s) = 2$. The machine's verdict, from a sieve of all ordered factorizations to one
million:

\begin{center}
\begin{tabular}{rrrr}
\hline
$y$ & $G(y)$ (boxes in the tower) & $C y^{\rho}$ & ratio \\
\hline
$10^3$ & $48{,}614$ & $4.882 \times 10^4$ & $0.9958$ \\
$10^4$ & $2{,}602{,}393$ & $2.614 \times 10^6$ & $0.9957$ \\
$10^5$ & $139{,}804{,}054$ & $1.399 \times 10^8$ & $0.9991$ \\
$10^6$ & $7{,}509{,}428{,}756$ & $7.491 \times 10^9$ & $1.0024$ \\
\hline
\end{tabular}
\end{center}

\noindent The measured slope over the last decade is $1.7301$ against $\rho = 1.7286$. Two features
of the table repay attention. First, the crystal genuinely is fractal: seven and a half billion boxes
march to a non-integer power law. Second, the ratio column does not converge monotonically; it
\emph{breathes} around $1$. The breathing is structure, not noise: the equation $\zeta(s) = 2$ has
non-real solutions as well, and in the theory of complex dimensions of self-similar structures
\cite{LapidusBook} such roots contribute log-periodic oscillations to the counting --- indeed
possessing non-real dimensions is that theory's \emph{definition} of fractality. Because the
crystal's similarity ratios $1/d$ have incommensurable logarithms, it is a non-lattice self-similar
structure in Lapidus's sense, and its complex dimensions are aperiodic. The same theory walks
directly back to the deepest open problem: by a theorem of Lapidus and Maier \cite{LapidusMaier}, the
inverse spectral problem for fractal strings --- ``can one hear the dimension?'' --- is solvable in
dimension $D$ exactly when $\zeta$ has no zeros on the line $\mathrm{Re}\, s = D$, so the Riemann
Hypothesis is equivalent to audibility in every dimension but $\tfrac12$.

The tower has a thermodynamic reading, and the reading has a name. Interpret $e^{-s}$ as a
Boltzmann factor: the free crystal is Julia's \emph{primon gas} \cite{Julia}, one bosonic mode per
axis at energy $\log p$, with partition function $\zeta(\beta)$ --- divergent at $\beta = 1$, the
free gas's Hagedorn point. The tower --- every ordered factorization a distinct microstate --- is
an interacting gas with partition function $G(\beta) = 1/(2 - \zeta(\beta))$, and its Hagedorn
point \emph{is the crystal's dimension}: divergence at $\zeta(\beta) = 2$, that is, at
$\beta = \rho = 1.7286\ldots$, driven by the $y^{\rho}$ density of states of Kalm\'ar's law ---
the tower gas boils while the free gas is still cold. The machine's check
(\texttt{primon\_gas.py}): thermal sums built from the exact factorization counts $K(n)$ to
$10^6$ match $1/(2-\zeta(\beta))$ with deficits equal to the $y^{\rho}$ density-law prediction to
ratios $0.999$--$1.000$ at every tested temperature, and the pole's residue
$1/(-\zeta'(\rho)) = 0.5500\ldots$ equals $\rho\, C$ with $C = 0.3181736\ldots$ the constant of
the table above: the counting law and the thermodynamics are one identity read twice. The
dimension of the crystal is the temperature at which its description tower boils.

A final calibration of how completely the complex plane reorganizes the crystal's bookkeeping. In the
real world, counting the layer-$2$ region to $y = 10^5$ means visiting $966{,}751$ lattice points one
by one. In the complex world the same information is stored at a single interior point: a numerical
contour integral around a circle of radius $0.4$ centred at $s = 1$ --- sampling $\zeta$ at a few
dozen complex values, nowhere near the lattice --- returns $966{,}735.679\ldots$, agreeing with the
theoretical residue to twenty decimal places and missing the true count by $15.3$ out of nearly a
million. The bulk lives at the pole; the leftover is the music.

\section{The diffraction pattern}

In the laboratory, hidden crystalline order announces itself by diffraction: shine waves on a
structure, and sharp Bragg peaks appear. Mathematical diffraction theory \cite{Hof, BaakeGrimm}
makes this rigorous for aperiodic structures: to a (possibly weighted) point comb one associates
its autocorrelation and takes the Fourier transform; the atoms of the result are the Bragg peaks.
One honesty clause from Section~2 governs everything below: the bare point sets $\{\log n\}$ and
$\{\log p\}$ are not Delone sets, so the geometric theory of model sets does not apply to them
directly. The objects that diffract are the \emph{weighted combs} the shadow carries, and for the
prime comb the diffraction statement is not an analogy but a theorem --- the Guinand--Weil
explicit formula \cite{Guinand, Weil52}. Paired against an even test function $h$ with
$g(u) = \frac{1}{2\pi}\int h(r)e^{-iur}\,dr$,
\[
\sum_{\gamma} h(\gamma) \;=\; \frac{1}{2\pi}\int_{-\infty}^{\infty} h(r)\,
\Bigl[\mathrm{Re}\,\psi\bigl(\tfrac14 + \tfrac{ir}{2}\bigr) - \log\pi\Bigr]dr
\;+\; h\bigl(\tfrac{i}{2}\bigr) + h\bigl(-\tfrac{i}{2}\bigr)
\;-\; 2\sum_{n \ge 2} \frac{\Lambda(n)}{\sqrt n}\, g(\log n),
\]
the left sum over all nontrivial-zero ordinates: the zero comb read at $h$ equals the
$\Lambda(n)/\sqrt n$-weighted log-prime comb read at $g$, plus the displayed pole and archimedean
terms. The identity is unconditional ($\gamma$ complex if zeros stray from the critical line;
under the Riemann Hypothesis the left side is the pairing of a genuine atomic measure on the real
line). The machine's check (\texttt{diffraction\_pairing.py}), with $h$ a Gaussian of unit width
centred at $\pm t_0$:

\begin{center}
\begin{tabular}{lccc}
\hline
centre $t_0$ & zero side & prime $+$ smooth side & $|$difference$|$ \\
\hline
$0$ (pole territory) & $0.000000000$ & $0.000000000$ & $1.3 \times 10^{-23}$ \\
$\gamma_1 = 14.1347\ldots$ & $2.000000000$ & $2.000000000$ & $7.6 \times 10^{-24}$ \\
$17.5$ (between zeros) & $0.010996589$ & $0.010996589$ & $2.8 \times 10^{-24}$ \\
$\gamma_2 = 21.0220\ldots$ & $2.000701572$ & $2.000701572$ & $3.9 \times 10^{-24}$ \\
$\gamma_3 = 25.0108\ldots$ & $2.000702435$ & $2.000702435$ & $3.5 \times 10^{-24}$ \\
$\gamma_4 = 30.4248\ldots$ & $2.085660945$ & $2.085660945$ & $4.3 \times 10^{-24}$ \\
\hline
\end{tabular}
\end{center}

\noindent At a zero the pairing is $\approx 2$, the $h$-mass of one comb atom counted with both
signs; between zeros it is small but nonzero and matched to twenty-three decimals all the same
--- an \emph{identity}, not a peak fit. At $t_0 = 0$ the pole term $h(\pm i/2) = 4e^{1/8}$ is
cancelled to the same precision by the archimedean and prime terms: the exact accounting behind
the ``bulk'' subtracted in the experiment below. Both experiments that follow are corollaries of
this display, read in its two directions.

\emph{Primes ring at zeros.} Take the prime powers up to $N = 2 \times 10^5$ at their logarithmic
positions, weighted as in the Chebyshev function, and scan the exponential sum
$\sum_{n \le N} \Lambda(n)\, n^{-1/2 - it}$ across frequencies $t$. The raw scan is dominated by the
smooth bulk term --- the pole's contribution, whose truncation ripple (spacing $2\pi/\log N$) is a
convincing impostor --- and the first lesson of the experiment is that the loudest structure in raw
data is usually the instrument. Subtract the bulk, and the true peaks emerge:

\begin{center}
\begin{tabular}{cccc}
\hline
peak located at $t$ & Riemann zero $\gamma$ & discrepancy \\
\hline
$14.140$ & $14.134725$ & $0.005$ \\
$21.010$ & $21.022040$ & $0.012$ \\
$25.035$ & $25.010858$ & $0.024$ \\
$30.380$ & $30.424876$ & $0.045$ \\
$32.950$ & $32.935062$ & $0.015$ \\
\hline
\end{tabular}
\end{center}

\noindent A sieve to $200{,}000$ and a wave scan recover the first five zeros of the zeta function to
within hundredths. The diffraction pattern of the primes is the zeros; and the raw scan's impostor
bulk was the pole term $h(\pm i/2)$ of the identity all along.

\emph{Zeros ring at primes.} The converse is a theorem of Landau from 1912 \cite{LandauZeros}: for
fixed $x > 1$, $\sum_{0 < \gamma \le T} x^{\rho}$ grows like $-(T/2\pi)\Lambda(x)$ --- it spikes
exactly when $x$ is a prime power, with amplitude proportional to $\log p$, and stays bounded
otherwise. With the first $300$ zeros ($T = 541.85$) and the detector normalized so that its limit is
$\Lambda(x)$, a scan of the integers $30$ to $60$ gives: every prime lit, at close to its true
$\log p$ (e.g.\ $53$ rang at $3.90$ against $\log 53 = 3.97$); the prime powers lit at the overtone
amplitudes ($49 = 7^2$ at $1.64$ against $\log 7 = 1.95$; $32 = 2^5$ faint but audible at $0.27$
against $\log 2 = 0.69$); and every composite silent, the loudest reaching $0.02$. The separation was
clean, with no exceptions. The prime powers' reduced amplitudes $\tfrac1k \log p$ are precisely the
``dust'' of the companion system \cite{CompanionI}, here made audible as the overtone series of the
prime tones.

For a true lattice, the statement ``the Fourier transform of the point set is again a point set'' is
Poisson summation. For the multiplication crystal's shadow, it is the Riemann--Weil explicit formula:
primes and zeros are each other's diffraction patterns, an observation Dyson made pointedly when he
proposed classifying quasicrystals as a road toward the Riemann Hypothesis \cite{Dyson}.

The modern vocabulary sharpens Dyson's road and prices it. A \emph{crystalline measure} is a
measure with locally finite support whose distributional Fourier transform is again one; a
\emph{Fourier quasicrystal} adds temperedness of the variations \cite{Meyer, LevOlevskii}. The
rigidity theorems of Lev and Olevskii \cite{LevOlevskii} say the tame positive examples ---
uniformly discrete support and spectrum --- are lattice combs in disguise; the Lee--Yang
constructions of Kurasov and Sarnak \cite{KurasovSarnak} produce genuinely exotic Fourier
quasicrystals, and were motivated by exactly this chapter's question: how special is the zero
comb? Where the universe stands, stated with labels. Theorem, unconditional: the explicit
formula displayed above. Theorems, cited: the rigidity and existence results
\cite{LevOlevskii, KurasovSarnak}. Conjecture (equivalent to RH): the zero comb
$\sum_\gamma \delta_\gamma$ is a measure on the \emph{real} line. Open: whether the zero comb ---
whose average spacing $2\pi/\log\gamma$ shrinks, so it is not uniformly discrete and escapes
every rigidity hypothesis --- admits a crystalline-measure characterization, and whether such a
characterization could be run backwards toward RH. The Kurasov--Sarnak examples arise as spectra
of quantum graphs --- metric graphs with incommensurable edge lengths, the same objects, with
lengths $\log 2, \log 3, \log 5, \ldots$, that this programme plans as its next instrument.

\section{The symmetries}

A crystallographer's first question about any crystal is its symmetry group, and here the crystal
delivers its deepest surprise: it is simultaneously the most symmetric object in elementary
mathematics and the least, and the distance between those two facts is a fair definition of number
theory.

\emph{Translations.} The crystal slides into itself under multiplication by every $d \ge 2$ --- but
only forward: division exits the positive integers. As it stands, the crystal is a semigroup crystal,
translatable one way; completing it to the positive rationals restores the full translation lattice
of infinite rank.

\emph{Maximal symmetry of the bare crystal.} Since $(\mathbb{N}, \times)$ is free on the primes, any
permutation of the primes extends to an automorphism, and these are all of them: the automorphism
group of the bare multiplication crystal is the full symmetric group on its axes. Multiplicatively,
nothing distinguishes $2$ from $3$: the map exchanging them inside every factorization preserves all
products. The bare crystal has no favorite directions at all.

\emph{Rigidity of arithmetic.} Now restore addition. The additive structure $(\mathbb{N}, +)$ is
generated by $1$, so its only automorphism is the identity; \emph{a fortiori} the automorphism group
of full arithmetic $(\mathbb{N}, +, \times)$ is trivial. Arithmetic is perfectly rigid. Welding a
maximally symmetric structure to a born-rigid one, along the enumeration of the primes by size, kills
every symmetry --- and the weld is the object of study. In the language of physics, the sizes
$\log 2, \log 3, \log 5, \ldots$ --- the incommensurable lengths that scrambled the shadow in
Section~2 --- are a \emph{symmetry-breaking field}, and the phenomena of the primes are the pattern
of the breaking. (The same field, incidentally, breaks music: $3^{12} \ne 2^{19}$ by the Pythagorean
comma of $1.4\%$, which is why no piano can be tuned perfectly. The untunability of pianos and the
irregularity of primes are one incommensurability.)

That language is not a metaphor on loan; it is a theorem with an address. Bost and Connes
\cite{BostConnes} built a quantum statistical-mechanical system out of exactly this chapter's
objects --- the crystal's translation semigroup acting on the congruence phases of the point
group --- whose partition function is $\zeta(\beta)$ and whose equilibrium (KMS) states undergo a
phase transition at $\beta = 1$: below it the equilibrium state is unique and symmetric; above
it, the extremal equilibrium states are permuted freely and transitively by the profinite point
group $\mathrm{Gal}(\mathbb{Q}^{\mathrm{ab}}/\mathbb{Q})$ of the clocks above --- spontaneous
symmetry breaking in the physicists' exact sense, with the symmetry that breaks being this
chapter's. The machine can watch the order parameter (\texttt{primon\_gas.py}): the Gibbs
expectation of the phase observable $e^{2\pi i/3}$, computed exactly through Hurwitz zeta
functions, is the frozen root of unity at $\beta = \infty$ (modulus $1.000000$, and at clock $5$
the four frozen states are the four primitive fifth roots, visibly permuted by the Galois group),
then falls as the system heats --- $0.893$ at $\beta = 4$, $0.529$ at $2$, $0.141$ at $1.2$,
$0.0076$ at $1.01$ --- condensing to zero at the transition: the congruence phases melt into the
unique symmetric state. The classification theorems are Bost's and Connes's, not ours; what the
crystal adds is the observation that their system is this chapter written as operator algebra ---
the symmetry that breaks is the point group, the field that breaks it is the spectrum of lengths,
and the transition sits at the pole of the free gas.

\emph{The logic of the weld.} The weld's wildness is itself a theorem --- three theorems, with
names. The bare ruler is decidable: Presburger proved in 1929 that the first-order theory of
$(\mathbb{N}, +)$ admits a decision procedure \cite{Presburger}. The bare crystal is decidable
too: Skolem proved it in 1930 for $(\mathbb{N}, \times)$ \cite{Skolem}, the modern proof running
axiswise --- a multiplicative question decomposes along the crystal's axes into ruler questions
about exponents. The weld is undecidable: G\"odel and Church \cite{Godel, Church} --- and Julia
Robinson sharpened the point to a knife: addition is definable from multiplication and
\emph{successor alone} \cite{JRobinson}, so the weld is fatal already at one tick of the ruler.
The machine walks all three (\texttt{logic\_column.py}): it decides Frobenius representability on
the ruler and brute force confirms the exact frontiers ($6x + 35y$ misses $169$ last; $20x + 27y$
misses $493$); it decides $\exists x, y\,(x^2 = 2y^2)$ --- the irrationality of $\sqrt{2}$ ---
negatively, mechanically, by the parity of a single exponent, a question the bare crystal can
settle but never ask about $x + 1$; and it verifies Robinson's identity
$a + b = c \iff S(ac)\,S(bc) = S(c^2\,S(ab))$ (for $c \neq 0$) at all $226{,}981$ triples of a
finite patch --- the definitional collapse in miniature. Two glosses. First, the logic column
separates the structures, not the universes: the mirror's weld $(\mathbb{F}_2[t], +, \times)$ is
undecidable too \cite{RRobinson}; the solvability the mirror enjoys elsewhere in this essay is
analytic, not logical, and the wildness everywhere is the weld's, not the lengths'. Second, the
laboratory of Section~8 has already leaned on this fact: the guarantee that no scanner exhausts
the weld's phenomena is exactly this undecidability, which here acquires its citations.

\emph{Surviving symmetry I: scale.} The crystal's natural measure $dx/x$ is exactly
dilation-invariant, and its spectrum under its own scaling symmetry --- the dilation frequencies that
resonate --- is the set of zeros: the diffraction experiments of Section~4 were representation theory
of the scale symmetry.

\emph{Surviving symmetry II: the congruence point groups.} On every clock $q$, the reduced residues
form the finite abelian group $(\mathbb{Z}/q)^{\times}$; the crystal's points distribute evenly over
it, and so, by Dirichlet's theorem \cite{Dirichlet}, do the primes. These finite groups play the role
of the point group, their characters --- waves on finite dials --- its irreducible representations;
all clocks assemble into the profinite point group at infinity, which by the Kronecker--Weber theorem
is precisely the symmetry group of the maximal abelian extension of the rationals \cite{Lang}. The
symmetry is exact only in the limit; at finite range it carries celebrated fine structure. Chebyshev
observed in $1853$ that the class $3 \bmod 4$ leads the race more often than not, a bias whose
logarithmic distribution Rubinstein and Sarnak determined, under natural hypotheses, on exactly the
phase torus the crystal hides \cite{RubinsteinSarnak}. And in $2016$ Lemke Oliver and Soundararajan
discovered, by counting, that \emph{consecutive} primes repel in their residues \cite{LOS}. Our own
census of the primes to $10^7$, by last digit of consecutive pairs:

\begin{center}
\begin{tabular}{c|cccc}
 & $\to 1$ & $\to 3$ & $\to 7$ & $\to 9$ \\
\hline
$1$ & $17.04\%$ & $31.08\%$ & $32.10\%$ & $19.77\%$ \\
$3$ & $23.16\%$ & $15.60\%$ & $29.27\%$ & $31.97\%$ \\
$7$ & $25.65\%$ & $27.55\%$ & $15.60\%$ & $31.20\%$ \\
$9$ & $34.12\%$ & $25.83\%$ & $23.07\%$ & $16.97\%$ \\
\end{tabular}
\end{center}

\noindent Same-digit pairs total $16.30\%$ against the naive $25\%$. The point-group symmetry is
exact for single primes and visibly broken for pairs: symmetry of the parts, fine structure of the
correlations.

\emph{Surviving symmetry III: the hidden Galois census.} Extending the number system recovers the
symmetries that rigidity forbids. Every prime acquires a symmetry label --- its Frobenius --- in the
Galois group of any extension, and Chebotarev's density theorem \cite{Lang} says the primes visit the
labels with perfect fairness: the primes are an exact census of the hidden symmetry group of
arithmetic. The harmonic analysis of that group is the Langlands program, which lies beyond this
essay; the point here is only where it sits on the crystallographic map: it is the representation
theory of the symmetries that survive extension.

\emph{Surviving symmetry IV: the mirror.} The reciprocal crystal --- the zeros --- possesses one
exact, proven symmetry: Riemann's functional equation \cite{Riemann}, $\xi(s) = \xi(1-s)$, a mirror
through the critical line. Machine check: at $s = 0.3 + 7.2i$,
$|\xi(s) - \xi(1-s)| = 3.6 \times 10^{-22}$ against $|\xi(s)| = 0.142$. In this chapter's language
the Riemann Hypothesis is a \emph{confinement} statement: the diffraction pattern, provably symmetric
about its mirror, is conjectured to lie \emph{on} it.

\emph{Surviving symmetry V: the statistical symmetry class.} Beyond exact symmetries, the reciprocal
crystal has a statistical one. Montgomery's pair-correlation discovery \cite{Montgomery}, confirmed
at scale by Odlyzko \cite{Odlyzko}, is that the zeros space themselves like the eigenvalues of large
random unitary matrices --- the GUE law, whose signature is quadratic repulsion of near-collisions.
Our test on the first $500$ zeros, spacings unfolded by the local density:

\begin{center}
\begin{tabular}{lccc}
\hline
spacing bin & observed & GUE & Poisson \\
\hline
$[0, 0.25)$ & $0.006$ & $0.016$ & $0.221$ \\
$[0.25, 0.5)$ & $0.062$ & $0.096$ & $0.172$ \\
$[0.5, 1)$ & $0.487$ & $0.421$ & $0.239$ \\
$[1, 1.5)$ & $0.343$ & $0.341$ & $0.145$ \\
$[1.5, 2)$ & $0.100$ & $0.108$ & $0.088$ \\
$[2, \infty)$ & $0.002$ & $0.017$ & $0.135$ \\
\hline
\end{tabular}
\end{center}

\noindent The smallest of $499$ normalized spacings is $0.236$; structureless (Poisson) points would
put about a hundred spacings below that value, and we observe one. The spectrum repels. In the
classification of Katz and Sarnak \cite{KatzSarnak} --- proven in the function-field universe,
conjectural in ours --- each zeta-like crystal carries a symmetry \emph{type} drawn from the
classical compact groups, and the multiplication crystal's type is unitary.

\section{The mirror: the crystal with quantized lengths}

Every difficulty met above condensed, on inspection, to one property: the crystal's axis lengths
$\log 2, \log 3, \log 5, \ldots$ share no common measure. There is a companion universe in which that
single property is switched off. In Weil's Rosetta stone --- the celebrated analogy, set out in his
1940 letter \cite{Krieger}, between number fields and function fields --- the role of the integers is
played by the monic polynomials over a finite field $\mathbb{F}_q$, the role of the primes by the
irreducible ones, and the crystal is built by the \emph{identical} free construction. The one change
is the size map: $|f| = q^{\deg f}$, so every axis length is an integer multiple of the quantum
$\log q$. The crystal becomes lattice rather than aperiodic, and the prediction of this essay's
thesis is stark: every aperiodic pathology of the integers should become an exactness. (In the
vocabulary of~\cite{Defs}: the table below evaluates one list of crystal invariants at two
crystals whose only difference is the multiset of axis lengths.) The machine's
verdict, chapter by chapter, over $\mathbb{F}_2[t]$:

\begin{center}\small
\begin{tabular}{p{0.20\textwidth} p{0.36\textwidth} p{0.36\textwidth}}
\hline
object & the integers & the mirror $\mathbb{F}_2[t]$ (verified) \\
\hline
prime count & no closed form; fluctuations tied to the zeros & exact:
$\tfrac1n \sum_{d \mid n} \mu(d)\, q^{n/d}$ (Gauss \cite{LidlNiederreiter}); verified to degree
$14$ \\
layer counts & the Piltz divisor problems; error exponents open since $1849$ & exact binomials
$q^n \binom{n-1}{k-1}$; verified for $k \le 6$, $n \le 12$ \\
layer birth & dimension $k$ opens at $2^k$ & at size $q^k$ (same law; both exact) \\
tower dimension & $\rho = 1.7286\ldots$, fractal; ratios breathe aperiodically & exactly
$1 + \log 2/\log q = 2$; tower mass $= 2^{2n-1}$ \emph{identically}, verified to $n = 14$ \\
complex dimensions & aperiodic (non-lattice self-similar) & exact periodic comb, period
$2\pi/\log q$ \\
M\"obius / Mertens & random walk; RH $\Leftrightarrow$ square-root size & identically zero in every
degree $\ge 2$; verified to degree $10$ \\
zeros & conjectured on the critical line & on $|\alpha| = \sqrt{q}$ by theorem
(Hasse--Weil \cite{Hasse, Weil49}); verified for $y^2 = x^3 + x + 1$ at all $428$ good primes
$\le 3000$, tightest approach $0.9917$ of the bound \\
spectral statistics & GUE, conjectural \cite{Montgomery, Odlyzko} & Sato--Tate, a theorem
\cite{SatoTate}; $428$ angles match the semicircle to $\sim 1\%$ per bin \\
rigidity & $\mathrm{Aut}(\mathbb{N}, +, \times) = 1$: arithmetic admits no symmetry & non-rigid:
$t \mapsto at + b$ are automorphisms; $t \mapsto t{+}1$ verified to permute the irreducibles of
every degree \\
point group & Dirichlet's theorem, with biases & effectively exact: degree $6$ splits
$3$--$3$--$3$ mod $t^2 + t + 1$ \\
the additive weld & the $abc$ conjecture, open & the Mason--Stothers theorem
\cite{MasonStothers} \\
logic of the weld & $(\mathbb{N},+)$ and $(\mathbb{N},\times)$ decidable (Presburger,
Skolem); the weld undecidable, already with successor & $(\mathbb{F}_2[t],+,\times)$
undecidable too (R. Robinson): the wildness is the weld's, not the lengths' \\
symmetry-breaking field & $\log p$ incommensurable & $\deg \cdot \log q$ quantized --- \emph{the measurable difference} \\
\hline
\end{tabular}
\end{center}

Two rows carry the section's meaning. The last row is the thesis made checkable: the two universes
share the crystal and differ in nothing but the arithmetic of the sizes, and every left-column
mystery sits opposite a right-column exactness. And the rigidity row explains \emph{why} the mirror
is solvable. The mirror's arithmetic possesses symmetries that ours provably lacks --- substitutions,
and above all the Frobenius --- and Weil's proof of the mirror's Riemann Hypothesis \cite{Weil49}
uses the Frobenius as the operator whose eigenvalues \emph{are} the zeros. The Hilbert--P\'olya dream
is not a dream there; it is the proof. The operator exists because the arithmetic has a symmetry for
it to be the operator of; our universe's rigidity is not a curiosity of the symmetry chapter but a
candidate explanation for a century of impasse. The function-field zeta function is Artin's
\cite{Artin}; the elliptic case of its Riemann Hypothesis is Hasse's \cite{Hasse}; the machine's role
here is only to watch the theorem work --- and to note, at one prime, the bound holding by less than
one percent: what our universe conjectures, the mirror enforces to the wire. A vertical verification completes
the picture: over $\mathbb{F}_5$ the same curve's two zeros
$\alpha, \bar\alpha = (-3 \pm i\sqrt{11})/2$ --- read off from a nine-point base count, and lying
on the circle $|\alpha| = \sqrt5$ by the exact identity $(9 + 11)/4 = 5$ --- predict the exact
number of points over every extension field, verified by brute force through $\mathbb{F}_{5^7}$
($78{,}633$ points among $78{,}125$ elements): seven exact predictions from two zeros, the Hasse
bound holding at every level, and the normalized traces equal to $|\cos n\theta|$ for the single
angle $\theta$ of the zero --- the mirror's explicit formula as one sustained tone, where ours
requires an infinite orchestra.

A third universe, checkable to the last digit: finite graphs carry the identical free
construction, with axes the primitive closed geodesics (backtrackless, tailless, up to rotation)
and lengths the integer edge-counts --- a lattice crystal with quantum $1$. Here the bookkeeping
is not merely exact but \emph{finite}: the crystal's zeta function, Ihara's \cite{Ihara}, is
rational by Bass's theorem \cite{Bass} --- for a $d$-regular graph,
$\zeta_G(u)^{-1} = (1-u^2)^{|E|-|V|}\det(I - Au + (d-1)u^2)$, a polynomial --- and its Riemann
Hypothesis (nontrivial poles on the circle $|u| = 1/\sqrt{d-1}$) is \emph{equivalent} to the
graph being Ramanujan: every nontrivial adjacency eigenvalue obeying $|\lambda| \le 2\sqrt{d-1}$
\cite{Terras, LPS}. The machine's verdict (\texttt{graph\_universe.py}; all graphs $3$-regular,
critical circle $|u| = 1/\sqrt2$): for $K_4$, $K_{3,3}$, the cube, and the Petersen graph, the
geodesic counts from the Bass determinant match the nonbacktracking traces integer for integer to
length $12$, the Euler product over primitive geodesics rebuilds the determinant exactly --- the
free monoid audited directly --- and every nontrivial pole sits on the critical circle to
$10^{-16}$; the circular ladder $CL_{16}$, whose spectral radius $2\cos(\pi/8) + 1 = 2.8478\ldots$
exceeds $2\sqrt2 = 2.8284\ldots$, fails Ramanujan, and its poles leave the circle in the same
breath ($8.8 \times 10^{-2}$ off), the equivalence holding in both directions. In this universe
the Riemann Hypothesis is not a conjecture one inherits but a property one measures, graph by
graph --- and the constructions of Lubotzky, Phillips, and Sarnak \cite{LPS} manufacture infinite
families that pass.

The ontological verdict this section was built to deliver: a prime is an irreducible element of a
free crystal, and \emph{that} is not mysterious --- the mirror dispatches it in closed form. A number
is a crystal point together with a size. The two universes differ only in the size map, and with it
goes everything: closed form versus none, integer dimension versus fractal breathing, periodic music
versus aperiodic, silence versus random walk, theorem versus conjecture. The mystery of the primes is
not the algebra of the atoms. It is the arithmetic of their lengths.

That verdict can be sharpened twice, and the sharpened form is this essay's final thesis. First, the
incommensurability has an \emph{address}. By Ostrowski's theorem \cite{Ostrowski}, every absolute
value on the rationals is either $p$-adic --- and every one of those is quantized, taking its values
in the powers of a single prime --- or it is the ordinary archimedean absolute value. Our universe is
therefore already quantized at all of its places but one, and the lengths $\log p$ are read, and
found incommensurable, precisely at that one place: the place at infinity, which an entire modern
field, Arakelov geometry, was built to tame. Second, quantization alone does not buy the mirror's
theorems: a Beurling system \cite{Beurling, DiamondZhang} with perfectly lattice lengths but
arbitrary multiplicities has periodic, meromorphic bookkeeping and, in general, neither a Riemann
Hypothesis nor a prime formula. The mirror's miracles require its second gift, the Frobenius --- the
symmetry $x \mapsto x^q$ that exists only in positive characteristic, powers the closed-form count,
and serves as the operator in Weil's proof --- and that engine keeps running: the twin-prime and
Chowla conjectures, frontier problems here, are theorems over $\mathbb{F}_q[t]$
\cite{SawinShusterman}, proved by the same mechanism of randomness derived from symmetry. And the
two gifts are one. Characteristic zero forces the archimedean place, hence the incommensurable
lengths; and characteristic zero forbids the Frobenius, hence the rigidity of the symmetry chapter
and the missing operator that Hilbert--P\'olya and its physical form \cite{BerryKeating} have
sought to retrofit. One dial, two absences, every mystery. The thesis, in its final compression: the
mystery of the primes is characteristic zero; the arithmetic of the lengths is its measurable
shadow.

\section{Between the universes: de-quantization along the temperament ladder}

The dictionary's last row treats commensurability as a switch. This section turns it into a dial. The
interpolation territory has two named maps. On the arithmetic side, Beurling's theory of generalized
primes \cite{Beurling, DiamondZhang} takes an arbitrary system of lengths, builds the free crystal on
them, and asks how much regularity of the lengths purchases which theorems; the answer is
quantitative and sharp --- a Prime Number Theorem is bought at error exponent $3/2$ in the integer
counts, and not below \cite{DMV}. On the dimension side, Lapidus's lattice/non-lattice dichotomy
\cite{LapidusBook} governs the complex dimensions: rational length-ratios force periodic combs,
irrational ones force aperiodic quasiperiodic scatter. Our experiment runs the two-length crystal
$\{2, 2^{p/q}\}$ up the continued-fraction convergents of $\log_2 3$ toward the true incommensurable
system $\{2, 3\}$ --- and the ladder's rungs come with names, because the denominators of those
convergents, $2, 5, 12, 41, 53, 306, 665$, are the historical equal temperaments of music:
twelve-tone equal temperament is the fourth rung, and the $41$- and $53$-tone systems of the
microtonalists are the fifth and sixth. The de-quantization path from the solvable universe to ours
runs through the pianos.

\begin{center}\small
\begin{tabular}{rrccc}
\hline
$p/q$ & temperament $q$ & dimension $D_{p/q}$ & comb lines & comb spacing $2\pi q/\ln 2$ \\
\hline
$3/2$ & $2$ & $0.811370$ & $3$ & $18.1$ \\
$8/5$ & $5$ & $0.783910$ & $8$ & $45.3$ \\
$19/12$ & $12$ & $0.788319$ & $19$ & $108.8$ \\
$65/41$ & $41$ & $0.787778$ & $65$ & $371.7$ \\
$84/53$ & $53$ & $0.787900$ & $84$ & $480.4$ \\
$485/306$ & $306$ & $0.787884$ & $485$ & $2773.8$ \\
$1054/665$ & $665$ & $0.787885$ & $1054$ & $6028.0$ \\
$\{2,3\}$ & (aperiodic) & $0.787885$ & --- & --- \\
\hline
\end{tabular}
\end{center}

Three findings. First, the dimensions oscillate onto the non-lattice limit exactly as convergents do
--- over, under, over --- agreeing with the measured dimension of the $\{2,3\}$ crystal to six
decimals by the $665$-tone rung. Second, the first rung is exactly solvable: the $\{2, 2^2\}$ tower
counts are the Fibonacci numbers, verified to satisfy the recurrence and Binet's rounding identity
for all $n \le 60$ --- lattice music as a single, perfectly periodic, geometrically damped overtone.
Third, the true system's complex dimensions do what the structure theory of \cite{LapidusBook} says
non-lattice strings must: a numerical scan to height $50$ finds nine roots (a count consistent with
the expected density $\ln 3/2\pi$), their real parts wandering across $[-0.988,\, 0.788]$ where every
rational rung pins them to vertical lines, the real dimension alone at the rightmost, and
quasi-periodic near-returns creeping back toward it ($+0.681$ at $t \approx 17.6$, $+0.767$ at
$t \approx 45.6$). Each finite rung is a comb whose period stretches with the temperament;
aperiodicity is the audible limit of ever-finer pianos.

The horizon, honestly drawn. This charts de-quantization for the self-similar skeleton of a
two-letter crystal; carrying the full integer crystal --- infinitely many incommensurable lengths,
welded to addition --- along such a path is the open frontier. At the far endpoint of the mirror,
where the quantum of length degenerates, live the field-with-one-element programs and the arithmetic
site of Connes and Consani \cite{ConnesConsani}, in which the crystal's translation semigroup becomes
a topos whose hidden points are conjectured to house the spectral realization of the zeros. The essay
ends at that door --- not through it, but knowing, with some precision, which door it is.

\section{The weld laboratory, the scanner, and the two-skeleton conjecture}

The essay so far has read the crystal's known structure. This section reports the experimental
programme built to ask whether anything \emph{else} is there, and states its organizing claim as a
conjecture equipped with its own falsification protocol.

\emph{The weld laboratory.} Selecting primes by multiplicative (crystal) criteria and then adding
them is a controlled experiment on the weld between the two arithmetics. Three modes were run. With
no selection, the representation counts of even numbers as sums of two primes display the celebrated
banded structure, and the machine measured its loudest band to four digits: multiples of $3$ receive
$1.999$ times the representations of their neighbours, against the Hardy--Littlewood singular
series' prediction of exactly $2$ \cite{HardyLittlewood}; dividing every count by the predicted
congruence factor collapses $25{,}001$ data points onto one curve with $2.26\%$ spread --- the entire
visible pattern of prime addition is the congruence skeleton in additive costume. A crystal-neighbour
mode (safe primes, $p = 2q+1$) shows the mechanism in miniature: the selection leaks a congruence ---
of the $4{,}324$ safe primes below $10^6$, all but the exception $7$ are $\equiv 2 \bmod 3$, a
three-line theorem --- and the leak casts an additive shadow: $99.95\%$ of pairwise sums are
$\equiv 1 \bmod 3$, and exactly one sum in all of mathematics is $\equiv 2 \bmod 3$, namely
$7 + 7 = 14$. A spectral mode (signs of $a_p$ on the Section~6 curve) returned the predicted null:
pair-sum label correlation within one standard error of independence. Recursing the neighbour mode
gives Cunningham chains, whose census below $10^6$ runs $67{,}134$, $5{,}953$, $949$, $109$, $24$,
$5$ by maximal length \cite{Dickson}, and whose starts obey a congruence \emph{amplifier}: at
recursion level $j$ each small prime $r$ forbids exactly one residue class mod $r$, the classes
distinct until the order of $2$ mod $r$ saturates, confining length-$4$ starts to $4$ of the $48$
coprime residues mod $210$ (machine-verified --- indeed the machine corrected the author's first
derivation of when the mod-$5$ filter binds). The dual recursion downward, through the factorizations
of $p-1$, is Pratt's primality certificate \cite{Pratt}, and it is shallow: mean depth $4.7$ and
maximum $10$ across all $68{,}906$ primes in $[10^5, 10^6]$, the doubly logarithmic height
established by Ford, Konyagin, and Luca \cite{FKL}.

\emph{The skeleton-null scanner.} To search without a theorem in hand, the programme uses a
pre-registered protocol: fix a finite battery of crystal-native invariants per prime (layer numbers
of $p \pm 1$, the Pratt height, the largest-factor slope, the unfolded gap, digit weight, characters)
and a fixed grid of same-prime and consecutive-prime correlations; take as null model not uniformity
but the \emph{skeleton} --- shuffles that preserve all congruence-class and density structure; pay
the multiple-comparison tax in advance (Bonferroni-grade thresholds over $117$ tests); demand
replication in an untouched holdout half; and audit every survivor up a ladder: instrument artifact,
mechanical feature entanglement, skeleton-derivable, known theorem, and only then candidate. First
light returned $69$ survivors, and the audit sorted the entire deck: a mechanical cluster (features
competing for one integer's size budget); a $2$-adic cluster that was the scanner \emph{finding its
own bug} (stratification by $p \bmod 30$ is blind to $p \bmod 4$, so the elementary seesaw between
$v_2(p-1)$ and $v_2(p+1)$ poured through the null); one detrending artifact (binary digit sums
against coarse size bands); and two genuine detections, both known --- the
Lemke Oliver--Soundararajan consecutive-residue repulsion \cite{LOS}, rediscovered unprompted at
$z = -27$ (the positive control passed), and the anticorrelation of consecutive gaps. Candidates:
zero, the honest harvest of any instrument's first light.

\emph{The extraction pipeline.} From scan to theorem the programme runs six stages: stabilize the
anomaly across scales; mechanize it with a skeleton-conditioned probabilistic model whose new
predictions are committed \emph{before} measurement; weaken the statement into the highest provable
stratum; test it in the mirror; prove and postdict the original readings; then the librarian and the
adversary. Executed live on the scanner's loudest reading (the $\chi_4$ cluster), the pipeline
delivered:

\begin{theorem}
For every prime $\ell$ and every $k \ge 1$, the density of primes $p$ with
$v_\ell(p-1) = k$ is exactly $\ell^{-k}$, with mass $1 - 1/(\ell - 1)$ at $k = 0$ and mean depth
$\ell/(\ell-1)^2$. In particular $\mathbb{E}[\,v_2(p-1) \mid p \equiv 3 \bmod 4\,] = 1$ identically
and $\mathbb{E}[\,v_2(p-1) \mid p \equiv 1 \bmod 4\,] \to 3$.
\end{theorem}

\noindent The proof is Dirichlet's theorem \cite{Dirichlet} read in the progressions
$1 \bmod \ell^k$ and telescoped; the factors $\ell - 1$ cancel. Machine verification: the $2$-adic
law to four decimals at $10^6$ and $10^7$; conditional means $2.9990$ and $1.0000$; the mirror law
over $\mathbb{F}_2[t]$ at mean $1.984$ against the predicted $2$. The pre-registration discipline
earned its keep: the committed prediction for the $3$-adic zero mass was $2/3$ (the random-integer
heuristic) and the measurement returned $0.5002$ --- Dirichlet's $1/2$ --- the programme's third
machine-corrects-the-author event, surfaced as a discrepancy rather than dissolved into hindsight.
The corollary (the class gap of $\mathbb{E}[\Omega(p-1)]$ tends to $2$, with the odd parts equalized
by Halberstam's shifted-prime Erd\H{o}s--Kac theorem \cite{Halberstam}) postdicts the scan's loudest
$z$ in sign and size. The librarian's verdict, applied without mercy: the theorem is classical
folklore. The pipeline's honest certificate for this input reads \emph{known result, freshly
mechanized, one operator error caught, one instrument reading explained} --- which is what
extraction usually yields, and should.

\emph{The conjecture.} The programme's organizing claim, assembled from every experiment above, is a
property of the crystal:

\begin{conjecture}[the two-skeleton conjecture]
The multiplication crystal's detectable order consists of its congruence skeleton and its spectral
skeleton, and nothing else. Operationally: for every finite pre-registered battery of crystal-native
invariants of the primes, every correlation detectable at any scale is derivable from the congruence
data together with the Hardy--Littlewood correlations; conditioned on these, all remaining
statistics agree with the skeleton-null ensemble. In particular the scanner of this section, run at
any scale with adequately resolved nulls, asymptotically returns no survivor beyond the
skeleton-derivable ones.
\end{conjecture}

\noindent Honesty about ownership and status. The mathematical content of this conjecture is not
new: its random side is the circle of conjectures of Chowla \cite{Chowla} and Sarnak \cite{Sarnak}
--- orthogonality of the M\"obius signal to all low-complexity observers --- and its structured side
is Hardy--Littlewood \cite{HardyLittlewood}; celebrated cases are theorems in the mirror
\cite{SawinShusterman}. The random side has proven base camps at low observer complexity ---
automatic sequences cannot correlate with M\"obius (M\"ullner \cite{Mullner}, after
Mauduit--Rivat's digits of primes \cite{MauduitRivat}), nor can bounded-depth circuits (Green
\cite{GreenAC0}), and the two-point logarithmic Chowla correlation vanishes (Tao \cite{TaoChowla})
--- so the conjecture is best read stratified by observer class, the scanner probing above the
theorem line; the ladder, the upgraded instrument, and its preregistered second first light are
the subject of the programme's Note~3 \cite{Ladder}. What the crystal formulation adds is operational: the conjecture names its
own experiment. Every audited null strengthens it; a single robust, replicated,
skeleton-inexplicable survivor refutes it; and by the undecidability of the weld it can never be
exhaustively verified --- the crystal guarantees its own inexhaustibility. The instrument, the
protocol, and the pipeline are in the repository \cite{Code}, standing equipment awaiting
candidates.

\section{What the crystal cannot tell us}

An essay conducted at this altitude owes the reader its ledger. Every load-bearing statement above is
a named theorem or a named conjecture: the geometric arithmetic is Euclid's; the dimension law is
Kalm\'ar's; complex dimensions and the audibility criterion are Lapidus's and Lapidus--Maier's; the
resonance of zeros at primes is Landau's; the explicit-formula duality is Riemann's and Weil's, its
quasicrystal reading Dyson's; equidistribution over the point group is Dirichlet's, over the Galois
labels Chebotarev's; the consecutive-prime bias is Lemke Oliver's and Soundararajan's; the race
analysis is Rubinstein's and Sarnak's; the mirror is Riemann's; the statistics are Montgomery's and
Odlyzko's, the symmetry types Katz's and Sarnak's; the primon gas is Julia's, and the phase
transition with its Galois symmetry is Bost's and Connes's. The mirror dictionary is Weil's
Rosetta stone \cite{Krieger}, its entries due to Gauss, Artin, Hasse, Weil, Stothers and Mason, and
the provers of Sato--Tate. The interpolation
territory of the final section is Beurling's and Lapidus's. The localization of the final thesis rests on Ostrowski's classification of the places and on the
function-field theorems of Sawin and Shusterman. The weld laboratory's asymptotic is Hardy and Littlewood's; the certificates are Pratt's; the
tree heights are Ford's, Konyagin's, and Luca's; the chain admissibility is Dickson's; the
shifted-prime input is Halberstam's; and the content of the two-skeleton conjecture is the
Chowla--Sarnak circle. What is offered as new is the organization --- the
crystal as a single object; Moran's reading of Kalm\'ar's equation; number theory as the physics of a
broken crystal symmetry --- together with the machine demonstrations, and the connection to the
companion papers, whose certified instruments \cite{CompanionII, CompanionIII} are, in this
language, \emph{certified crystallography}: exact, provable measurements of the point group and its
breaking, as opposed to the statistical observations tabulated here.

Two epistemic warnings, both earned in our own computations. First, serene data prove nothing on this
terrain: the strict form of square-root cancellation suggested by every computation ever performed
(Mertens's conjecture) is false, disproved through the zeros without any explicit counterexample
being known. Second, the modern structure-versus-randomness programme sharpens the crystal's
limits: its detectable order appears to consist of exactly the two skeletons measured in this essay
--- the congruence skeleton and the spectral skeleton --- and the conjectures of Chowla and Sarnak
assert that beyond them no low-complexity detector finds anything at all \cite{Sarnak}. The preceding
section states this as the two-skeleton conjecture and equips it with a falsification protocol. The crystal
shows us its symmetries and their breaking; what remains after both are subtracted is, on the field's
considered bet, genuine noise.

\section{Reproducibility}

Every table in this essay is produced by a single documented script (\texttt{crystal\_demos.py}) in
the code repository \cite{Code}, using only integer sieves, the mpmath and numpy libraries, and, for
Section~4, the first $300$--$500$ zeros of $\zeta$ computed on the spot. Total runtime is a few minutes on ordinary hardware; the mirror computations (irreducible sieves
over $\mathbb{F}_2[t]$, the elliptic-curve point counts, and the Sato--Tate histogram) add under a
minute. The graph-universe checks of Section~6 (\texttt{graph\_universe.py}), the
explicit-formula pairing table of Section~4 (\texttt{diffraction\_pairing.py}), and the
primon-gas and Bost--Connes order-parameter tables of Sections~3 and~5
(\texttt{primon\_gas.py}), and the decision-procedure demonstrations of Section~5
(\texttt{logic\_column.py}) add seconds each. The weld laboratory, the recursion census, the skeleton-null scanner, and the extraction run add
a few minutes more. The companion papers' exact systems \cite{CompanionI, CompanionII,
CompanionIII} were re-used where indicated and were themselves machine-verified as described there.

\section*{Coda}

The multiplication crystal is maximally symmetric; arithmetic is perfectly rigid; and everything we
call the mystery of the primes lives between those two sentences, in the breaking of the one symmetry
by the sizes of the axes. The primes are the census of the symmetry that hides; the zeros are the
spectrum of the symmetry that survives; and the oldest unsolved problem in mathematics is the
conjecture that the spectrum keeps to its mirror. Compressed once more: the
crystal is universal; characteristic zero is ours --- one dial, two absences, every mystery. A reader who takes away one image should take the
plane of ones: everything in this essay --- the fractal tower, the Bragg peaks, the broken point
group --- was already latent in the child's picture of numbers as pebbles arranged in rows, looked at
long enough, and counted honestly.

\section*{Disclosure of AI use}

In accordance with publisher policies on the use of artificial intelligence, the author discloses the
following. This essay was developed in an extensive collaboration with Claude (Anthropic; Claude
Fable 5, accessed through the claude.ai interface, 2026). The direction of the work --- the founding
intuition that arithmetic should be read as the geometry of a space of ones, pursued through a
sequence of the author's conjectures about projection, dimension, and symmetry --- is the author's,
together with all decisions about scope, claims, and submission. Within that direction, the AI system
identified the classical results corresponding to those conjectures and their attributions, wrote and
ran the verification code producing every numerical table, and drafted and revised the text; the
reason for its use was to enable a single non-specialist author to develop, check, and document the
material to a publishable standard. The author has reviewed the entire manuscript and takes full
responsibility for its content, including the accuracy of the mathematics and of all references.

\begin{thebibliography}{99}
\bibitem{CompanionI} C. Gribble, Prime counts and prime sums from multiplication tables: a
self-similar formula system, companion manuscript, Zenodo (2026), \texttt{doi:10.5281/zenodo.21769103}.
\bibitem{CompanionII} C. Gribble, Certified exact sums of primes over dyadic intervals via a ladder
of generalized factorials, companion manuscript, Zenodo (2026), \texttt{doi:10.5281/zenodo.21769105}.
\bibitem{CompanionIII} C. Gribble, Recovering the primes in a dyadic interval from exact moments: two
constructive routes and an obstruction theorem, companion manuscript, Zenodo (2026),
\texttt{doi:10.5281/zenodo.21769107}.
\bibitem{Defs} C. Gribble, Crystals and their shadows: definitions for a programme on the
multiplication crystal, programme note (2026), in the code repository
\texttt{https://github.com/carlgribble-caa/prime-moments}.
\bibitem{Euclid} T. L. Heath, \emph{The Thirteen Books of Euclid's Elements}, Dover, New York, 1956.
\bibitem{HardyRamanujan} G. H. Hardy and S. Ramanujan, The normal number of prime factors of a number
$n$, \emph{Quart. J. Math.} \textbf{48} (1917), 76--92.
\bibitem{Kalmar} L. Kalm\'ar, \"Uber die mittlere Anzahl der Produktdarstellungen der Zahlen,
\emph{Acta Litt. Sci. Szeged} \textbf{5} (1931), 95--107.
\bibitem{LapidusBook} M. L. Lapidus and M. van Frankenhuijsen, \emph{Fractal Geometry, Complex
Dimensions and Zeta Functions}, 2nd ed., Springer, New York, 2013.
\bibitem{LapidusMaier} M. L. Lapidus and H. Maier, The Riemann hypothesis and inverse spectral
problems for fractal strings, \emph{J. London Math. Soc.} (2) \textbf{52} (1995), 15--34.
\bibitem{LandauZeros} E. Landau, \"Uber die Nullstellen der Zetafunktion, \emph{Math. Ann.}
\textbf{71} (1912), 548--564.
\bibitem{Dyson} F. Dyson, Birds and frogs, \emph{Notices Amer. Math. Soc.} \textbf{56} (2009),
212--223.
\bibitem{Hof} A. Hof, On diffraction by aperiodic structures, \emph{Comm. Math. Phys.}
\textbf{169} (1995), 25--43.
\bibitem{BaakeGrimm} M. Baake and U. Grimm, \emph{Aperiodic Order. Vol. 1: A Mathematical
Invitation}, Cambridge Univ. Press, 2013.
\bibitem{Guinand} A.~P. Guinand, A summation formula in the theory of prime numbers, \emph{Proc.
London Math. Soc.} (2) \textbf{50} (1948), 107--119.
\bibitem{Weil52} A. Weil, Sur les ``formules explicites'' de la th\'eorie des nombres premiers,
\emph{Comm. S\'em. Math. Univ. Lund} (1952), 252--265.
\bibitem{Meyer} Y. Meyer, Measures with locally finite support and spectrum, \emph{Proc. Natl.
Acad. Sci. USA} \textbf{113} (2016), 3152--3158.
\bibitem{LevOlevskii} N. Lev and A. Olevskii, Quasicrystals and Poisson's summation formula,
\emph{Invent. Math.} \textbf{200} (2015), 585--606.
\bibitem{KurasovSarnak} P. Kurasov and P. Sarnak, Stable polynomials and crystalline measures,
\emph{J. Math. Phys.} \textbf{61} (2020), 083501.
\bibitem{Riemann} B. Riemann, \"Uber die Anzahl der Primzahlen unter einer gegebenen Gr\"osse,
\emph{Monatsber. K\"onigl. Preuss. Akad. Wiss. Berlin} (1859), 671--680.
\bibitem{Dirichlet} P. G. L. Dirichlet, Beweis des Satzes, dass jede unbegrenzte arithmetische
Progression \ldots unendlich viele Primzahlen enth\"alt, \emph{Abhandl. K\"onigl. Preuss. Akad.
Wiss.} (1837), 45--81.
\bibitem{Lang} S. Lang, \emph{Algebraic Number Theory}, 2nd ed., Springer, New York, 1994.
\bibitem{RubinsteinSarnak} M. Rubinstein and P. Sarnak, Chebyshev's bias, \emph{Experiment. Math.}
\textbf{3} (1994), 173--197.
\bibitem{LOS} R. J. Lemke Oliver and K. Soundararajan, Unexpected biases in the distribution of
consecutive primes, \emph{Proc. Natl. Acad. Sci. USA} \textbf{113} (2016), E4446--E4454.
\bibitem{Montgomery} H. L. Montgomery, The pair correlation of zeros of the zeta function, in
\emph{Analytic Number Theory}, Proc. Sympos. Pure Math. XXIV, Amer. Math. Soc., 1973, 181--193.
\bibitem{Odlyzko} A. M. Odlyzko, On the distribution of spacings between zeros of the zeta function,
\emph{Math. Comp.} \textbf{48} (1987), 273--308.
\bibitem{KatzSarnak} N. M. Katz and P. Sarnak, Zeroes of zeta functions and symmetry, \emph{Bull.
Amer. Math. Soc.} \textbf{36} (1999), 1--26.
\bibitem{Julia} B. Julia, Statistical theory of numbers, in \emph{Number Theory and Physics},
Springer Proc. Phys. 47, Springer, 1990, 276--293.
\bibitem{BostConnes} J.-B. Bost and A. Connes, Hecke algebras, type III factors and phase
transitions with spontaneous symmetry breaking in number theory, \emph{Selecta Math.} \textbf{1}
(1995), 411--457.
\bibitem{Presburger} M. Presburger, \"Uber die Vollst\"andigkeit eines gewissen Systems der
Arithmetik ganzer Zahlen, in welchem die Addition als einzige Operation hervortritt, in
\emph{Comptes Rendus du I Congr\`es des Math\'ematiciens des Pays Slaves}, Warsaw, 1929, 92--101.
\bibitem{Skolem} T. Skolem, \"Uber einige Satzfunktionen in der Arithmetik, \emph{Skrifter utgitt
av Det Norske Videnskaps-Akademi i Oslo} I, no.~7 (1930).
\bibitem{Godel} K. G\"odel, \"Uber formal unentscheidbare S\"atze der Principia Mathematica und
verwandter Systeme I, \emph{Monatsh. Math. Phys.} \textbf{38} (1931), 173--198.
\bibitem{Church} A. Church, An unsolvable problem of elementary number theory, \emph{Amer. J.
Math.} \textbf{58} (1936), 345--363.
\bibitem{JRobinson} J. Robinson, Definability and decision problems in arithmetic, \emph{J.
Symbolic Logic} \textbf{14} (1949), 98--114.
\bibitem{RRobinson} R.~M. Robinson, Undecidable rings, \emph{Trans. Amer. Math. Soc.}
\textbf{70} (1951), 137--159.
\bibitem{Sarnak} P. Sarnak, Three lectures on the M\"obius function, randomness and dynamics, lecture
notes, Institute for Advanced Study, 2011.
\bibitem{MauduitRivat} C. Mauduit and J. Rivat, Sur un probl\`eme de Gelfond: la somme des chiffres
des nombres premiers, \emph{Ann. of Math.} \textbf{171} (2010), 1591--1646.
\bibitem{Mullner} C. M\"ullner, Automatic sequences fulfill the Sarnak conjecture, \emph{Duke Math.
J.} \textbf{166} (2017), 3219--3290.
\bibitem{GreenAC0} B. Green, On (not) computing the M\"obius function using bounded depth circuits,
\emph{Combin. Probab. Comput.} \textbf{21} (2012), 942--951.
\bibitem{TaoChowla} T. Tao, The logarithmically averaged Chowla and Elliott conjectures for
two-point correlations, \emph{Forum Math. Pi} \textbf{4} (2016), e8.
\bibitem{Ladder} C. Gribble, The observer ladder: stratifying the two-skeleton conjecture, with a
second first light, programme note (2026), in the code repository
\texttt{https://github.com/carlgribble-caa/prime-moments}.
\bibitem{Krieger} M. H. Krieger, A 1940 letter of Andr\'e Weil on analogy in mathematics,
\emph{Notices Amer. Math. Soc.} \textbf{52} (2005), 334--341.
\bibitem{LidlNiederreiter} R. Lidl and H. Niederreiter, \emph{Finite Fields}, 2nd ed., Cambridge
Univ. Press, 1997.
\bibitem{Artin} E. Artin, Quadratische K\"orper im Gebiete der h\"oheren Kongruenzen I, II,
\emph{Math. Z.} \textbf{19} (1924), 153--246.
\bibitem{Hasse} H. Hasse, Zur Theorie der abstrakten elliptischen Funktionenk\"orper III, \emph{J.
Reine Angew. Math.} \textbf{175} (1936), 193--208.
\bibitem{Weil49} A. Weil, Numbers of solutions of equations in finite fields, \emph{Bull. Amer.
Math. Soc.} \textbf{55} (1949), 497--508.
\bibitem{MasonStothers} W. W. Stothers, Polynomial identities and hauptmoduln, \emph{Quart. J.
Math. Oxford} \textbf{32} (1981), 289--305; R. C. Mason, \emph{Diophantine Equations over Function
Fields}, Cambridge Univ. Press, 1984.
\bibitem{SatoTate} T. Barnet-Lamb, D. Geraghty, M. Harris, and R. Taylor, A family of Calabi--Yau
varieties and potential automorphy II, \emph{Publ. Res. Inst. Math. Sci.} \textbf{47} (2011),
29--98.
\bibitem{Ihara} Y. Ihara, On discrete subgroups of the two by two projective linear group over
$\mathfrak{p}$-adic fields, \emph{J. Math. Soc. Japan} \textbf{18} (1966), 219--235.
\bibitem{Bass} H. Bass, The Ihara--Selberg zeta function of a tree lattice, \emph{Internat. J.
Math.} \textbf{3} (1992), 717--797.
\bibitem{Terras} A. Terras, \emph{Zeta Functions of Graphs: A Stroll through the Garden},
Cambridge Univ. Press, 2011.
\bibitem{LPS} A. Lubotzky, R. Phillips, and P. Sarnak, Ramanujan graphs, \emph{Combinatorica}
\textbf{8} (1988), 261--277.
\bibitem{Beurling} A. Beurling, Analyse de la loi asymptotique de la distribution des nombres
premiers g\'en\'eralis\'es I, \emph{Acta Math.} \textbf{68} (1937), 255--291.
\bibitem{DiamondZhang} H. G. Diamond and W.-B. Zhang, \emph{Beurling Generalized Numbers}, Math.
Surveys Monogr. 213, Amer. Math. Soc., Providence, 2016.
\bibitem{DMV} H. G. Diamond, H. L. Montgomery, and U. M. A. Vorhauer, Beurling primes with large
oscillation, \emph{Math. Ann.} \textbf{334} (2006), 1--36.
\bibitem{ConnesConsani} A. Connes and C. Consani, The arithmetic site, \emph{C. R. Math. Acad. Sci.
Paris} \textbf{352} (2014), 971--975.
\bibitem{Ostrowski} A. Ostrowski, \"Uber einige L\"osungen der Funktionalgleichung
$\varphi(x)\varphi(y) = \varphi(xy)$, \emph{Acta Math.} \textbf{41} (1916), 271--284.
\bibitem{SawinShusterman} W. Sawin and M. Shusterman, On the Chowla and twin primes conjectures over
$\mathbb{F}_q[T]$, \emph{Ann. of Math.} \textbf{196} (2022), 457--506.
\bibitem{BerryKeating} M. V. Berry and J. P. Keating, The Riemann zeros and eigenvalue asymptotics,
\emph{SIAM Rev.} \textbf{41} (1999), 236--266.
\bibitem{HardyLittlewood} G. H. Hardy and J. E. Littlewood, Some problems of `Partitio
numerorum' III: On the expression of a number as a sum of primes, \emph{Acta Math.} \textbf{44}
(1923), 1--70.
\bibitem{Chowla} S. Chowla, \emph{The Riemann Hypothesis and Hilbert's Tenth Problem}, Gordon and
Breach, New York, 1965.
\bibitem{Pratt} V. R. Pratt, Every prime has a succinct certificate, \emph{SIAM J. Comput.}
\textbf{4} (1975), 214--220.
\bibitem{FKL} K. Ford, S. V. Konyagin, and F. Luca, Prime chains and Pratt trees, \emph{Geom.
Funct. Anal.} \textbf{20} (2010), 1231--1258.
\bibitem{Halberstam} H. Halberstam, On the distribution of additive number-theoretic functions III,
\emph{J. London Math. Soc.} \textbf{31} (1956), 14--27.
\bibitem{Dickson} L. E. Dickson, A new extension of Dirichlet's theorem on prime numbers,
\emph{Messenger of Math.} \textbf{33} (1904), 155--161.
\bibitem{Code} C. Gribble, Reference implementations accompanying this essay and its companions, code
repository (2026), \texttt{https://github.com/carlgribble-caa/prime-moments}.
\end{thebibliography}

\end{document}
